\section{Discussion}
\label{sec:discussion}

% =====================================================================
% §7.1 — What the null means for policy
% =====================================================================
\subsection{What the null means for deinstitutionalization}
\label{sec:discussion-policy}

Deinstitutionalization has been litigated for decades on a prediction:
that closing psychiatric hospitals faster than communities can absorb
the displaced patients raises mortality, and suicide in particular,
among the most vulnerable. Brazil implemented the reform at national
scale. On the populations with realized patient-flow ties to the closing
hospitals, rather than the merely nearby, we find no evidence of a large
acute mortality increase. Population-weighted suicide mortality among
the exposed changed by $\valSuicWAtt$ per $100{,}000$
(95\% CI $[\valSuicWLo, \valSuicWHi]$); self-harm was likewise null; parallel trends
hold on the backfilled data.

The policy reading must respect the bound. The evidence excludes an
increase larger than about a quarter of baseline suicide among the
exposed, which would lie outside the confidence interval. It cannot
exclude a smaller effect, which $\valNclosuresPnash$ closures do not
resolve. The finding is therefore not a blanket safety claim. It is that,
in the Brazilian case, the data do not show the large acute mortality
increase its critics feared in the populations that bore the closures.
For a reform implemented worldwide largely without
population-scale mortality evidence, an identified null of that
magnitude, on the right population, is informative on its own terms. It
also shifts where the policy debate should look: if acute mortality did
not spike, the welfare question turns to the slower-moving margins of
chronic disease management, family burden, and quality of life, which
our mortality outcomes do not capture.

This interpretation is complementary to \citet{diasFontes2024}. Their
CAPS rollout estimates show substitution away from psychiatric
hospitalizations, no reduction in mental-health mortality, and higher
violent crime. Our closure-level estimates ask whether the hospital
side of the same reform produced acute mortality increases among the
municipalities that had actually used the closing facilities.

% =====================================================================
% §7.2 — The measurement lesson
% =====================================================================
\subsection{The measurement lesson for closure studies}
\label{sec:discussion-generalize}

The null is only as credible as the population it is estimated on, and
that is where the paper's portable contribution lies. The closure
literature defines exposure by geographic proximity
\citep{avdic2016, carroll2019, joynt2015, lindrooth2018}; in our setting
that rule misclassifies $\valNmisclass$ of the $\valNunionPairs$ municipality--closure pairs,
missing the patients who crossed borders for specialized care and
flagging neighbors who never used the facility. A distance-based
analysis of the same closures would have estimated effects on a
substantially different, partly disjoint population---and would have
diluted the flow-dependent municipalities with the merely nearby. The
patient-flow exposure rule, checked by a mechanical first stage (travel
burden falls where patients had realized utilization ties to the closing
hospital), addresses this measurement problem.

The correction is not specific to Brazil or to psychiatry. It requires
three conditions: specialized care concentrated in regional hubs so that
proximity and dependence diverge; origin--destination claims data from a
single or near-universal payer; and a setting where patients cross
administrative borders at scale. These hold across many single-payer and
single-payer-like systems. Where they hold, closure-impact studies that
define exposure by generic distance may estimate effects on a different
population than the one that used the closing hospital. The size and
direction of the discrepancy are measurable ex ante from the flow data,
as our divergence map shows.

% =====================================================================
% §7.3 — Why the cross-section was uninformative
% =====================================================================
\subsection{Why we identify only from within-event variation}
\label{sec:discussion-crosssection}

A natural alternative would compare mortality across municipalities by
their flow isolation, without a closure shock. That cross-sectional
comparison is uninformative here, for reasons that also explain why the
within-event design is necessary. Vital-registration coverage is
systematically worse in the rural North and Northeast---precisely the
regions with the largest flow-versus-distance access gaps---so any
cross-sectional access--mortality gradient is attenuated by measurement
quality that correlates with the regressor. And cross-municipal
comparisons confound flow isolation with the entire regional
health-system architecture. Both problems are absorbed by municipality
fixed effects under a closure shock; neither is in a cross-section. We
accordingly rely only on within-event variation, and we note (consistent
with the broader access-measurement literature, e.g.\
\citealp{doyle2011}) that cross-sectional spending- or
access-on-mortality regressions can fail for exactly these reasons. The
eleven cross-sectional specifications, which confirm the embedding adds
no out-of-sample predictive content over distance, are in
Appendix~\ref{app:xs}.

% =====================================================================
% §7.4 — Limitations
% =====================================================================
\subsection{Limitations}
\label{sec:discussion-limits}

Four limitations bound the claims. First, the closure sample is modest
($\valNclosuresPnash$ primary, $\valNclosuresPsymax$ at the widest), so the null rules out large effects
but not small ones, and fine heterogeneity is underpowered; the result
does, though, survive Federation-Unit clustering and randomization
inference (Section~\ref{sec:robustness}), so the modest sample does not
manufacture the null through understated standard errors. Second, the
analysis observes the SUS-using population only; private inpatient flows
are unrecorded, and whether the flow-versus-distance misalignment is
similar in high-private-penetration regions is open. Third, mortality
captures the acute margin; the slower welfare consequences of
deinstitutionalization---chronic disease trajectories, caregiver burden,
quality of life---lie outside it. Fourth, the closure-timing argument
rests on the PNASH inspection calendar and the reimbursement freeze,
which are institution-level and plausibly orthogonal to local demand
conditional on the design checks but are not randomized; the
natural-language motive classification that documents closure narratives
has internal but not yet external validation (Appendix~\ref{app:nlp}), so we use it only for
robustness sub-samples. None of these overturns the identified result;
each marks where it stops.
